\chapter{Návrh architektúry systému}
\label{NavrhArchitektury}

Táto kapitola popisuje celkovú architektúru navrhovaného riadiaceho systému. Vychádza z~požiadaviek definovaných v~predchádzajúcej kapitole a~z~analýzy dostupných technológií. Popisuje hardvérovú topológiu, softvérovú štruktúru, komunikačnú vrstvu a~formát konfiguračných súborov.

% ------------------------------------------------------------------
\section{Požiadavky na systém}
\label{PoziadavkyNaSystem}

Táto kapitola definuje hranice projektu a~podmienky, ktoré musí výsledný riadiaci systém spĺňať pre zabezpečenie spoľahlivej prevádzky v~prostredí interaktívnej expozície.

Z~hľadiska \textbf{funkčných požiadaviek} sa očakáva najmä bezproblémová integrácia a~synchronizácia rôznorodých hardvérových prvkov:
\begin{itemize}
    \item \textbf{Riadenie motora a~výkonu:} Primárnou úlohou je bezdrôtové riadenie výkonových prvkov, špecificky jednosmerných motorov s~asynchrónnym riadením ich rýchlosti a~smeru otáčania. Ďalšou kľúčovou funkciou je priame spínanie záťaží pracujúcich na sieťovom napätí \textbf{230\,V} (osvetlenie, dymostroje) a \textbf{12\,V}.
    \item \textbf{Multimediálna synchronizácia:} Keďže moderné expozície silne stavajú na audiovizuálnych vnemoch, systém musí zabezpečiť natívne a~synchronizované prehrávanie zvuku a~videa priamo z~vlastného úložiska.
    \item \textbf{Interaktivita a~časová os (\textit{timeline}):} Systém musí vedieť na základe hardvérových vstupov lokačne asynchrónne zmeniť vetvenie prebiehajúceho deja, no zároveň podporovať striktné chronologické časovanie sekvencií.
\end{itemize}

Medzi \textbf{nefunkčné požiadavky} s~ohľadom na udržateľnosť patria:
\begin{itemize}
    \item \textbf{Modularita a~softvérové oddelenie:} Pre zachovanie oddelenia softvérového inžinierstva od obsahovej tvorby sa úprava logiky (scenára) musí dať vykonať výhradne modifikáciou \textbf{textových konfiguračných súborov}. Tým sa odstraňuje nutnosť rekompilácie C++ firmware-u alebo zásahu do zdrojových kódov Pythonu.
    \item \textbf{Dostupnosť (Cost-effectiveness):} Vzhľadom na finančné limity inštalácií musí byť systém postavený na komerčne dostupných, cenovo výhodných mikrokontroléroch a~\textbf{Internet of Things (IoT)} platformách namiesto extrémne drahých dedikovaných priemyselných automatov (PLC).
\end{itemize}

% ------------------------------------------------------------------
\section{Logická a~sieťová koncepcia}
\label{CelkovaArchitektura}

Na architektúru navrhovaného systému je nevyhnutné nazerať z~dvoch odlišných, no vzájomne sa dopĺňajúcich hľadísk: z~\textbf{logického (výpočtového)} a~z~\textbf{fyzického (sieťového)}. 

Z~logického hľadiska ide o~\textbf{distribuovanú architektúru}. Výpočtový výkon a~inteligencia nie sú sústredené výhradne v~jednom centrálnom procesore (ako je bežné pri starších centralizovaných PLC automatoch), ale sú rozložené do priestoru priamo k~výkonovým prvkom. Koncové periférie postavené na mikrokontroléroch neslúžia len ako pasívne spínače. Naopak, aktívne a~autonómne dekódujú prichádzajúce textové správy (parsovanie reťazcov MQTT payloadu), prekladajú ich na strojové požiadavky a~priamo generujú špecifické signály (napr. pulzno-šírkovú moduláciu pre plynulé riadenie otáčok motorov).

Z~fyzického a~sieťového hľadiska však systém prísne dodržiava \textbf{hviezdicovú topológiu (\textit{Star topology})}. Hoci sú distribuované uzly plne inteligentné, z~dôvodu udržania kontroly a~eliminácie cyklenia na sieti nekomunikujú vzájomne medzi sebou (netvoria \textit{Mesh} sieť). Akýkoľvek dátový tok na sieti vedie striktne obojsmerne len cez centrálny uzol, ktorým je Wi-Fi prístupový bod so sprevádzkovaným MQTT sprostredkovateľom.

\vspace{1cm}
\begin{center}
    \textit{[MIESTO PRE OBRÁZOK: Bloková schéma fyzickej topológie hviezdicovej siete s~odlíšením centrálneho uzla a okrajových distribuovaných periférií]}
\end{center}
\vspace{1cm}

% ------------------------------------------------------------------
\section{Hardvérová voľba a~sieťový prenos}
\label{HardverovaTopoligia}

Systém si z~titulu svojej distribuovanej logiky s~hviezdicovou topológiou vyžiadal rozdelenie hardvéru do dvoch samostatných kategórií: na jeden centrálny uzol a~na sieť výkonných periférií. Z~hľadiska spoľahlivosti a~rýchlosti prenosu medzi nimi by bolo ideálne zabezpečiť komunikáciu výhradne plne káblovým (drôtovým) spojením. V~praxi však v~expozíciách či \textit{escape room} inštaláciách často nie je technologicky alebo esteticky možné naťahovať dodatočné káblové zväzky naprieč miestnosťou. Z~tohto dôvodu bolo na plošnú distribúciu dátových príkazov z~centrály zvolené bezdrôtové Wi-Fi pripojenie.

Dvojúrovňová architektúra bola fyzicky obsadená nasledovným hardvérom:
\begin{itemize}
    \item \textbf{Centrálny uzol (\uv{Mozog}):} O~zastávanie funkcie hlavného režiséra inštalácie sa stará jednodoskový počítač \textbf{Raspberry Pi}. Tento model bol zvolený z~niekoľkých praktických dôvodov. Okrem toho, že je tu sprevádzkovaný hlavný logický automat, zachovanie si celkového prehľadu nad výstavou kombinuje  s~priamym vyhradením výkonu na prehrávanie pamäťovo náročných médií vďaka natívnym HDMI a~analógovým audio výstupom. Zároveň dokáže súbežne hostovať ovládací webový \textit{dashboard} pre obsluhujúci personál, čím je riadenie fyzicky separované od exponátov.
    \item \textbf{Distribuované uzly (\uv{Svaly a zmysly}):} Na plnenie lokálnych fyzických úkonov slúži modulárna sieť bezdrôtových mikrokontrolérov \textbf{ESP32}, ktoré sú implementované a~napájané priamo pri interaktívnych objektoch. Tieto uzly sa na základe svojho zamerania softvérovo profilujú na:
    \begin{itemize}
        \item \textbf{Senzorický uzol} (tzv. \textit{Trigger Button}) pre hardvérovú detekciu vstupov od používateľov.
        \item \textbf{Motorický uzol} (\texttt{controller\_MOTORS}) pre obsluhu fyzického pohybu exponátov (jednosmerné motory).
        \item \textbf{Reléový uzol} (\texttt{controller\_RELAY}) pre spínanie efektov a~osvetlenia či iných zariadení pracujúcich na sieťovom napätí 230\,V.
    \end{itemize}
\end{itemize}

% ------------------------------------------------------------------
\section{Princíp softvérového jadra (Scene Engine)}
\label{SoftverovArchitektura}

Návrh riadiacej softvérovej logiky si vyžiadal zásadný odklon od bežného sekvenčného (zhora-nadol) programovania. Lineárne vykonanie bloku so zaradením vyčkávacích prvkov (napr. umelé uspatie vlákna) spôsobuje stratu kontroly nad asynchrónnymi udalosťami (napríklad ojedinelé stlačenie dekompresného tlačidla návštevníkom počas prehrávania zvuku). Z~tohto dôvodu bolo softvérové jadro architektované ako neblokujúci \textbf{Konečný stavový automat (FSM -- \textit{Finite State Machine})}.

Priebeh expozície je z~architektonického hľadiska logicky rozdelený na \textbf{samostatné oddelené stavy}. Životný cyklus každého stavu je riadený vlastnými inicializačnými, bežiacimi a~ukončovacími blokmi, ktoré na seba nečakajú v~zmysle funkčného blokovania, ale asynchrónne reagujú na časové značky a~vonkajšie udalosti. 

Systém neustále na pozadí vyhodnocuje \textit{udalosti z~prostredia} (napr. vypršanie maximálneho času v~stave, koniec prehrávania média, alebo príchod konkrétnej hardvérovej správy). Akonáhle je splnená prechodová podmienka, \textbf{mechanizmus dynamických prechodov} bezpečne preruší bežiaci stav a~stroj sa plynulo pretransponuje do nového cieľového stavu, čím vo výsledku vzniká interaktívne vetvenie výstavy na základe vôle návštevníka. Konkrétna softvérová implementácia týchto blokov v~jazyku Python a formáte JSON bude detailne rozobraná v~Kapitole \ref{ImplementaciaSW}.

\vspace{1cm}
\begin{center}
    \textit{[MIESTO PRE OBRÁZOK: Diagram konečného stavového automatu (FSM) -- prechody z onEnter do timeline a čakanie na transition eventy]}
\end{center}
\vspace{1cm}

% ------------------------------------------------------------------
\section{Architektúra komunikačnej vrstvy}
\label{KomunikacnaVrstva}

Pre okamžitú asynchrónnu výmenu dát medzi centrálnym automatom (v~jazyku Python) a~distribuovanými perifériami (ESP32) bez nadbytočnej réžie bol na komunikáciu zvolený \textbf{protokol MQTT} (\textit{Message Queuing Telemetry Transport}). Prenos informácií prebieha na osvedčenom koncepte \textit{Publish/Subscribe} cez konštantné TCP/IP pripojenie na centrálny sieťový sprostredkovateľ (\textit{broker}).

Navrhnutý model siete implementuje niekoľko kľúčových komunikačných štandardov:
\begin{itemize}
    \item \textbf{Štruktúra zverejňovania (\textit{topic} mriežka):} Architektúra prísne hierarchicky oddeľuje výstavy a~koncové prvky. Koncové ESP32 mikrokontroléry odoberajú výhradne svoj priradený topik (napríklad \texttt{room1/motor1}). 
    \item \textbf{Spätná väzba (\textit{feedback}):} Pri aktivácii príkazu centrálny systém zverejní povel (\textit{payload}) a~následne sa prepína do režimu vyčkávania. Periférny uzol potvrdí reálne fyzické vykonanie príkazu odpovedaním na identický topik obohatený o~príponu \texttt{/feedback}.
    \item \textbf{Detekcia dostupnosti cez LWT (\textit{Last Will and Testament}):} Sledovanie \textit{online/offline} stavu uzlov úmyselne nerieši aktívne dotazovanie siete (\textit{pingovanie}). Opiera sa o~zdieľanú zostatkovú (\textit{retained}) správu udržiavanú \textit{brokerom}, ktorej okamžitý prepis pri strate spojenia s uzlom efektívne zalarmuje centrálu o~poruche.
    \item \textbf{Dôsledná separácia dát z~kódu (JSON formát):} Vrcholom systémovej architektúry je abstrahovanie logiky. Kompletná choreografia udalostí \textit{timeline}, počiatočných akcií \textit{onEnter} i~dynamických prechodov je vyčlenená mimo tvrdý kód. Tieto dáta sú načítavané priamo z~externých konfiguračných súborov \texttt{.json}.
\end{itemize}

Vďaka textovému kontraktu JSON sa editorom podoby a logiky výstavy môže stať priamo samotný používateľ vo svojom textovom editore, a~to v~režime chodu programu bez najmenšej potreby zasahovania či rekompilácie jadra aplikácie.

\vspace{1cm}
\begin{center}
    \textit{[MIESTO PRE OBRÁZOK: Schéma toku MQTT správ (Publish/Subscribe) medzi Raspberry Pi a koncovými ESP32 uzlami so spätnou väzbou]}
\end{center}
\vspace{1cm}
